\section{Stage 1 results: data}\label{sec:stage1-results}

This section is filled automatically. After the steps of \cref{sec:stage1}
have been run on the selected video, \cmd{memvid export-tex} writes the tables
below from the manifest; until then, placeholders are shown. The numbers are
never typed by hand.

\subsection{Selected video}
\textit{Describe the selected video here: title, source and licence, genre,
duration, why it was chosen (recurring entities, state changes, speech), and
any candidates that were rejected and why.}

\subsection{Preprocessing}
\generated{stage1_results.tex}{preprocessing statistics, produced by
\cmd{memvid preprocess} and \cmd{memvid export-tex}.}

\paragraph{Checklist for interpreting the table.}
\begin{itemize}[nosep]
  \item \emph{Speech coverage} and the \emph{longest gap without speech} show
        how much of the video carries information only visually; they bound
        how useful speech-only facts can be.
  \item \emph{Suspicious segments} are candidates for hallucinated text; any
        that reach the memory should be discussed.
  \item The \emph{maximum frame timing error} must be below one frame period;
        a larger value indicates a decoding problem.
\end{itemize}
